\documentclass[12pt,a4paper]{report}
\usepackage[utf8]{inputenc}
\usepackage{booktabs}
\usepackage{geometry}
\usepackage{hyperref}
\usepackage{titlesec}
\usepackage{longtable}
\usepackage{xcolor}
\usepackage{enumitem}
\geometry{margin=1in}

\titleformat{\chapter}[hang]{\normalfont\huge\bfseries}{\thechapter.}{20pt}{\Huge}

\title{\textbf{Graph Databases --- Main Path Project (PathX)}}
\author{University of Padua \\[0.5em]
\textbf{Course:} Graph Databases \\
\textbf{Professor:} Gianmaria Silvello \\[1em]
\textbf{Team:} Juliette Tonneau, Minseok Kim, Bita Khashechian, Danil Smirnov}
\date{Academic Year 2025/26}

\begin{document}

\maketitle
\tableofcontents
\newpage

\chapter{Project Description}
This repository contains the implementation of the \textbf{Main Path} project for the \textit{Graph Databases} course.

The goal is to build a \textbf{system that integrates an SQL database} with \textbf{a Large Language Model (LLM)}, enabling natural language query processing and query optimization inspired by the \textit{Galois} paper presented in the course.

\chapter{Stage 1 — Tech Selection and System Setup}
This stage focuses on:
\begin{itemize}
    \item Selecting and setting up the development environment,
    \item Preparing a local SQL database (DuckDB),
    \item Implementing a Python pipeline that:
    \begin{enumerate}
        \item Connects to the database,
        \item Runs SQL queries from a \texttt{.sql} file,
        \item Saves results in a \texttt{.json} output file.
    \end{enumerate}
\end{itemize}

\noindent The teacher must be able to execute a single Python command to generate a JSON with result sets for each query. (Reference: Table 2 of the paper.)

\chapter{GALOIS Paper Overview and Key Insights}
The SIGMOD 2025 paper \textit{"Logical and Physical Optimizations for SQL Query Execution over Large Language Models"} presents Galois, a system that enables SQL query execution directly over Large Language Models (LLMs).

Galois introduces innovations that redefine query optimization to balance execution cost and result quality, introducing operators such as table-scan and key-scan. It dynamically builds metadata and estimates confidence scores based on model self-assessment.

Key lessons for implementation:
\begin{itemize}
    \item Introduce logical and physical plans that adapt query decomposition to the LLM's strengths and limitations
    \item Treat the LLM as a data source, not just a reasoning tool
    \item Maintain a cost-quality balance
    \item Structure results in JSON to ensure compatibility with analysis tools
\end{itemize}

\chapter{Implementation Schedule}

\section*{Phase 1 – Tech Selection and System Setup}
\begin{longtable}{lll}
\toprule
\textbf{Deadline} & \textbf{Task} & \textbf{Description} \\
\midrule
18 Oct & Literature review & Study and summarize the Galois SIGMOD 2025 paper. \\
19–20 Oct & System architecture design & Define system components and architecture diagram. \\
21 Oct & Python environment setup & Create virtual environment, install libraries. \\
21 Oct & Database setup & Install DuckDB/PostgreSQL, import datasets. \\
24 Oct & SQL connection test & Write Python scripts to connect, query, and retrieve JSON results. \\
25 Oct & EXPLAIN plan test & Execute \texttt{EXPLAIN} queries for logical plan analysis. \\
30–31 Oct & Final verification & Run end-to-end tests and commit setup. \\
\bottomrule
\end{longtable}

\section*{Phase 2 – NL and SQL Baselines}
\begin{longtable}{lll}
\toprule
\textbf{Date} & \textbf{Task} & \textbf{Description} \\
\midrule
1–3 Nov & NL baseline scripting & Write prompts for natural-language queries. \\
4–6 Nov & SQL baseline scripting & Write equivalent SQL queries. \\
7–8 Nov & LLM orchestration & Integrate baselines into LangChain workflow. \\
9–11 Nov & Prompt strategy testing & Experiment with prompt variations. \\
12–14 Nov & Evaluation metrics setup & Implement evaluation library (accuracy, recall, F1). \\
15–17 Nov & Unit testing & Build test runner for JSON output validation. \\
18–19 Nov & Documentation & Prepare README + PDF. \\
\bottomrule
\end{longtable}

\section*{Phase 3 – Implementation of Galois Features (Part I)}
\begin{longtable}{lll}
\toprule
\textbf{Date} & \textbf{Task} & \textbf{Description} \\
\midrule
20–22 Nov & Implement LLMScan Operator & Develop class to query LLM with table name and attributes. \\
23–26 Nov & Implement Filter-LLMScan Operator & Extend to handle WHERE clauses. \\
27–29 Nov & Multi-relation joining & Add joins and tuple merging. \\
30 Nov–2 Dec & Two-phase retrieval & Implement key-first (Key-Scan) retrieval. \\
3–4 Dec & End-to-end test & Run EU Cities example and collect metrics. \\
5 Dec & Demo preparation & Prepare demo comparing approaches. \\
\bottomrule
\end{longtable}

\section*{Phase 4 – Implementation of Galois Features (Part II)}
\begin{longtable}{lll}
\toprule
\textbf{Date} & \textbf{Task} & \textbf{Description} \\
\midrule
6–8 Dec & Galois Without Optimizations & Implement baseline version for Table 3 replication. \\
9–11 Dec & Testing and metrics & Run experiments and compute metrics. \\
12 Dec & Advisor demo & Present informal results. \\
13–16 Dec & Implement Galois Selective & Add logic for selective conditions. \\
17–19 Dec & Implement Galois All Attribute & Extend optimization to all attributes. \\
20–22 Dec & Implement Galois Full & Combine optimizations and finalize model. \\
23 Dec & Final verification & Replicate Table 3 and finalize demo. \\
\bottomrule
\end{longtable}

\section*{Phase 5 – Final Presentation}
\begin{longtable}{lll}
\toprule
\textbf{Date} & \textbf{Task} & \textbf{Description} \\
\midrule
2–4 Jan & Slide preparation & Draft slides summarizing results. \\
5–7 Jan & Rehearsal & Practice and finalize slides. \\
8 Jan & Final submission & Submit slides and present. \\
\bottomrule
\end{longtable}

\chapter{Technology Stack}
\begin{longtable}{lll}
\toprule
\textbf{Component} & \textbf{Tool / Library} & \textbf{Description} \\
\midrule
Language & Python 3.11+ & Main programming language \\
Database & DuckDB & Lightweight SQL database \\
LLM (later) & IBM WatsonX.ai / OpenAI API & Large Language Model integration \\
Libraries & duckdb, pandas, sqlglot, python-dotenv & Core dependencies \\
IDE & Visual Studio Code & Development environment \\
\bottomrule
\end{longtable}

\chapter{Installation and Setup}
\section*{Clone the repository}
\begin{verbatim}
git clone https://github.com/Graph-Databases-2025-26/PathX.git
cd PathX
\end{verbatim}

\section*{Create a virtual environment}
\begin{verbatim}
python -m venv .venv
source .venv/bin/activate     # macOS / Linux
.venv\Scripts\activate        # Windows
\end{verbatim}

\section*{Install dependencies}
\begin{verbatim}
pip install -r requirements.txt
\end{verbatim}

\section*{Initialize DuckDB and load data}
\begin{verbatim}
python scripts/init_duckdb.py
\end{verbatim}

\section*{Typical usage}
\begin{verbatim}
python -m venv .venv
source .venv/bin/activate
pip install -r requirements.txt
python scripts/build_ground_truth.py --data-root data --ground-root ground_truth_result \
  --datasets flight-2 flight-4 fortune geo premier presidents world movies

python scripts/galois_eval.py --ground ground_truth_result --submissions ground_truth_result \
  --datasets GEO MOVIES PRESIDENTS WORLD FLIGHT-2 FLIGHT-4 --cell-metric similarity --tuple-metric constraint
\end{verbatim}

\end{document}
